\documentclass[twocolumn]{iopjournal}

\usepackage{amsmath,amssymb,amsfonts}

\fancyhfoffset[L]{0mm}
\setlength{\marginparwidth}{0pt}
\setlength{\marginparsep}{0pt}

% ==========================================
% ZENODO PREPRINT HEADERS (ACTIVE)
% ==========================================
\renewcommand{\articletype}[1]{%
  {\vspace*{-8mm}\noindent \Large \sf Preprint}\par
  \vspace*{6mm}\noindent{\scriptsize \sf{\bfseries \MakeUppercase{#1}}}}

\fancyhead[L]{}
\fancyhead[R]{}

\hypersetup{pdftitle={Chiral Holonomy of the Closed Tetrad Frame}, pdfauthor={Stephen Nelson}, pdfsubject={Gravity Field Dynamics (GFD)}, pdfkeywords={holonomy, parallel transport, stellated octahedron, cube, chirality, spinor, double cover, proof}}

\begin{document}

\articletype{Supplementary Material}

\title{Chiral Holonomy of the\\
Closed Tetrad Frame}

\author{Stephen Nelson}

\affil{Field Dynamics Research}

\ead{stephen@fielddynamics.org}

\begin{abstract}
The eight null directions of the closed tetrad frame are the
corners of a cube, four on TetA and four on TetB, and every cube edge
joins the two. We transport a frame around the cube's Hamiltonian
cycles in two ways: \emph{uniformly}, with the rotation $R_{ij}$ that
carries $v_i$ to $v_j$ at every step, and \emph{chirally}, with its
transpose $R_{ij}^T$ at every step leaving a TetB direction.
Transposition keeps the diagonal and exchanges the two triangles.
Uniform transport is parallel transport on the sphere. Every
Hamiltonian cycle encloses exactly half the sphere, so the holonomy is
the identity on vectors and exactly $-1$ on spinors: a full turn,
undone only by a second circuit. Chiral transport never closes. Its
holonomy angles are exactly $\arccos(-4639/6561)$ and
$\arccos(-5983/6561)$, one for each sense of circulation, so the
closure is chiral. Exchanging the labels TetA and TetB is the same as
reversing the sense. Weighting each step by the shares $\mu$ and
$1-\mu$ gives a rotation by $\arctan(\sigma\tan\theta)$ with
$\sigma = (x-1)/(x+1)$; the chirality is odd in $\sigma$ and vanishes
at $x = 1$.
\end{abstract}

\keywords{holonomy, parallel transport, stellated octahedron, cube, chirality, spinor, double cover, proof}

% ----------------------------------------------------------------------

\section*{Setting}

Let $v \in \{(\pm1, \pm1, \pm1)/\sqrt3\}$ be the eight corners of the
cube on the unit sphere, the null directions of the closed
frame~\cite{L1,L3holonomy}. Following~\cite{L3holonomy}, TetA is the
set with $v_1v_2v_3 > 0$ and TetB the set with $v_1v_2v_3 < 0$; parity
$v \to -v$ exchanges them. (Theorem~3 shows the choice of label does
not matter.) Corners $v_i, v_j$ are adjacent when they
differ in one coordinate; then they lie in opposite sets and
$v_i\cdot v_j = 1/3$, so the angle between them is
$\theta = \arccos(1/3)$. The cube graph has six Hamiltonian cycles,
twelve directed cycles from any given corner.

For adjacent $i, j$ let $n_{ij} = v_i\times v_j/|v_i\times v_j|$, $K$
its cross-product matrix, and
\begin{equation}
  R_{ij} = I + \sin\theta\,K + (1 - \cos\theta)\,K^2 ,
  \label{eq:R}
\end{equation}
the rotation by $\theta$ about $n_{ij}$. For a directed cycle
$C = (c_0, \dots, c_7)$ define
\begin{equation}
\begin{aligned}
  H_u(C) &= R_{c_7c_0}\cdots R_{c_1c_2}R_{c_0c_1},\\
  H_\chi(C) &= T_{c_7c_0}\cdots T_{c_1c_2}T_{c_0c_1},
\end{aligned}
  \label{eq:H}
\end{equation}
with $T_{ij} = R_{ij}$ if $v_i \in$ TetA and $T_{ij} = R_{ij}^T$ if
$v_i \in$ TetB. We call these the uniform and the chiral transport.

\medskip
\noindent
\textit{Reading.} In the tetrad connection each $K\times K$ slice is
$6 + 4 + 6$: an upper triangle, a shared diagonal and a lower triangle,
$64 = 24 + 16 + 24$ in all, with TetA and TetB on the two
triangles~\cite{Gram}. Transposition is exactly the exchange of the
triangles with the diagonal kept (Lemma~1). The torsion-free
connection keeps $16 + 24 = 40$ and cannot tell the two roots
apart~\cite{L2}; uniform transport, reading every step as TetA, is its
analogue. These readings motivate the definitions; no proof below
uses them.

% ----------------------------------------------------------------------

\section*{Two Lemmas}

\paragraph{Lemma 1 (the edge rotation).} \textit{$R_{ij}$ is a
rotation with $R_{ij}v_i = v_j$, and
$R_{ij}^T = R(n_{ij}, -\theta) = R_{ji} = R_{ij}^{-1}$. Transposition
leaves the diagonal of $R_{ij}$ unchanged and exchanges its strict
upper and lower triangles.}

\paragraph{Proof.} $v_i\perp n_{ij}$ and Eq.~\eqref{eq:R} turns the
plane $\perp n_{ij}$ by $\theta$, from $v_i$ towards $v_j$.
$K^T = -K$ and $(K^2)^T = K^2$, which give $R_{ij}^T$; and
$n_{ji} = -n_{ij}$ gives $R_{ji}$. $K$ has zero diagonal and
$K^2$ is symmetric. \hfill $\square$

\paragraph{Lemma 2 (turning and area).} \textit{Along a Hamiltonian
cycle drawn with great-circle arcs, every turn is $\pm\pi/3$, and
every Hamiltonian cycle makes four turns of each sign. The region on
its left therefore has area exactly $2\pi$.}

\paragraph{Proof.} The three edges at a corner leave it at
$120^\circ$ to one another (by the threefold symmetry about that
corner), so a path entering on one and leaving on another turns by
$\pm\pi/3$. The count of four and four is checked on all six cycles.
By Gauss--Bonnet on the unit sphere the region on the left has area
$2\pi - \sum(\text{turns}) = 2\pi$. \hfill $\square$

% ----------------------------------------------------------------------

\section*{Uniform Transport}

\paragraph{Theorem 1.} \textit{For every Hamiltonian cycle,
$H_u(C) = I$, while the $SU(2)$ lift of the transport, with step
$U_{ij} = \cos\tfrac\theta2\,\mathbf 1 - i\sin\tfrac\theta2\,
n_{ij}\cdot\boldsymbol\sigma$, is exactly $-\mathbf 1$. Two circuits
give $+\mathbf 1$. Reversing the cycle transposes $H_u$.}

\paragraph{Proof.} $R_{ij}$ is the Levi-Civita parallel transport of
the sphere along the arc from $v_i$ to $v_j$: it is the rotation about
the arc's pole, which carries the tangent plane along the arc. Around
a closed geodesic polygon the transport is the rotation about the
starting point by the enclosed area, here $2\pi$ (Lemma~2), hence
$H_u = I$. For spin $\tfrac12$ the same loop gives the factor
$e^{-i\Omega/2}$ with solid angle $\Omega = 2\pi$~\cite{Berry}, that is
$-1$. Both are confirmed by exact computation on all twelve cycles,
each $U_{ij}$ covering $R_{ij}$. Reversal replaces each $R_{ij}$ by
$R_{ji} = R_{ij}^T$ in the opposite order, which is $H_u^T$.
\hfill $\square$

\medskip
\noindent
One circuit is a full turn: vectors return, spinors return with a
minus sign, and only a second circuit restores them. Uniform
transport cannot tell the two senses of circulation apart, since
$H_u$ and $H_u^T$ are the same rotation.

% ----------------------------------------------------------------------

\section*{Chiral Transport}

\paragraph{Theorem 2.} \textit{For every Hamiltonian cycle,
$H_\chi(C) \neq I$, and its rotation angle $\Theta$ satisfies}
\begin{equation}
  \cos\Theta \in \left\{-\frac{4639}{6561},\;
  -\frac{5983}{6561}\right\}, \qquad 6561 = 3^8 ,
  \label{eq:angles}
\end{equation}
\textit{that is $\Theta = 134.996^\circ$ or $155.770^\circ$. The two
senses of each undirected cycle give the two values. The $SU(2)$
half-traces are $31/81$ and $-17/81$. $\Theta$ does not depend on the
starting corner.}

\paragraph{Proof.} Exact computation on all twelve cycles, with
$\cos\theta = 1/3$, $\sin\theta = 2\sqrt2/3$, $\cos\tfrac\theta2 =
\sqrt{2/3}$. The $SU(2)$ values agree with Eq.~\eqref{eq:angles}
through $\cos\Theta = 2\cos^2\tfrac\Theta2 - 1$:
$2(31/81)^2 - 1 = -4639/6561$ and $2(17/81)^2 - 1 = -5983/6561$.
Moving the starting corner conjugates $H_\chi$. \hfill $\square$

\medskip
\noindent
The first angle is close to $3\pi/4$ but not equal to it:
$\cos(3\pi/4) = -0.70711$ against $-4639/6561 = -0.70706$.

\paragraph{Theorem 3 (relabelling is reversal).} \textit{Let
$H'_\chi$ be the chiral transport with the labels TetA and TetB
exchanged. Then, for every cycle,}
\begin{equation}
  H'_\chi(C) = H_\chi(\bar C)^T ,
  \label{eq:relabel}
\end{equation}
\textit{where $\bar C$ is $C$ traversed in the opposite sense. So
exchanging the two chambers exchanges the two values of
Eq.~\eqref{eq:angles}, exactly as reversing the sense does.}

\paragraph{Proof.} Write $M_k = T_{c_kc_{k+1}}$, so
$H_\chi(C) = M_7\cdots M_0$. The step of $\bar C$ from $c_{k+1}$ to
$c_k$ leaves the other chamber. By Lemma~1 it is
$R_{c_{k+1}c_k} = R_{c_kc_{k+1}}^T$ if $c_{k+1} \in$ TetA, that is,
if $c_k \in$ TetB, and $R_{c_{k+1}c_k}^T = R_{c_kc_{k+1}}$ otherwise.
Either way it equals $M_k$, so $H_\chi(\bar C) = M_0M_1\cdots M_7$.
Relabelling transposes every step, so
$H'_\chi(C) = M_7^T\cdots M_0^T = (M_0\cdots M_7)^T$.
\hfill $\square$

\medskip
\noindent
The two senses use the same eight steps in opposite orders. For
uniform transport the orders give transposes, the same rotation; for
chiral transport they give different rotations. That difference is
the chirality.

% ----------------------------------------------------------------------

\section*{Weighting by the Shares}

Weight each step by the shares $\mu$ and $1-\mu$ of the constitutive
law, $\mu(x) = x/(x+1)$~\cite{Exchange}: $S_{ij} = \mu R_{ij} +
(1-\mu)R_{ij}^T$ leaving TetA and $\mu R_{ij}^T + (1-\mu)R_{ij}$
leaving TetB. Let $H_\sigma(C)$ be the product, with
$\sigma = 2\mu - 1$.

\paragraph{Proposition 4.} \textit{(a) $S_{ij} = I +
\sigma\sin\theta\,K + (1-\cos\theta)K^2$ leaving TetA. It is normal,
with eigenvalue $1$ along $n_{ij}$ and $\cos\theta \pm
i\sigma\sin\theta$ in the plane: a rotation by $\varphi$, with}
\begin{equation}
  \tan\varphi = \sigma\tan\theta ,
  \label{eq:phi}
\end{equation}
\textit{times a contraction by $\rho = (\cos^2\theta +
\sigma^2\sin^2\theta)^{1/2} \ge 1/3$. (b) $\sigma = (x-1)/(x+1)$, the
difference of the shares, and $x \to 1/x$ sends $\sigma \to -\sigma$.
(c) $H_{-\sigma}(C) = H_\sigma(\bar C)^T$. (d) At $\sigma = 0$
($x = 1$) every step is symmetric and the two senses have the same
rotational part; the difference between the senses is odd in
$\sigma$. $\sigma = \pm1$ give Theorem~2.}

\paragraph{Proof.} (a) Substitute Eq.~\eqref{eq:R}; $S_{ij}$ is a
polynomial in $K$, hence normal, and the plane block is
$\begin{pmatrix}\cos\theta & -\sigma\sin\theta\\ \sigma\sin\theta &
\cos\theta\end{pmatrix}$. (b) $2x/(x+1) - 1$. (c) Leaving TetB the
step at $\sigma$ is the TetA step at $-\sigma$, its transpose; the
argument of Theorem~3 applies. (d) At $\sigma = 0$ the step is
$I + (1-\cos\theta)K^2$, symmetric, and by (c)
$H_0(C) = H_0(\bar C)^T$; a matrix and its transpose have
transposed polar factors, hence equal rotation angles. By (c) the
difference changes sign with $\sigma$. \hfill $\square$

% ----------------------------------------------------------------------

\section*{Summary}

On the cube of null directions, uniform transport is parallel
transport on the sphere. Each Hamiltonian cycle encloses half the
sphere, so it returns vectors unchanged and spinors with a minus sign.
Chiral transport, with the TetB steps transposed, never closes. Its
two exact angles belong to the two senses of circulation, and
exchanging the chambers exchanges them. Weighted by the shares of the
constitutive law, each step rotates by $\arctan(\sigma\tan\theta)$,
and the chirality is odd in $\sigma = (x-1)/(x+1)$ and vanishes at
$x = 1$.

% ----------------------------------------------------------------------

\begin{thebibliography}{9}

\bibitem{L3holonomy}
S.~Nelson,
\textit{Supplementary material: chiral holonomy of the closed tetrad
frame}, Letter~3 supplement, Field Dynamics Research (2026).

\bibitem{L1}
S.~Nelson,
\textit{The coordinate-space geometry of the tetrad frame},
Field Dynamics Research (2026).

\bibitem{Gram}
S.~Nelson,
\textit{The Gram matrix of the tetrad equation and the stellated
octahedron}, Field Dynamics Research (2026).

\bibitem{L2}
S.~Nelson,
\textit{The other half of gravity: torsion and the two-chamber
structure of the tetrad frame}, Field Dynamics Research (2026).

\bibitem{Exchange}
S.~Nelson,
\textit{The constitutive law from exchange symmetry},
Field Dynamics Research (2026).

\bibitem{Berry}
M.~V.~Berry,
\textit{Quantal phase factors accompanying adiabatic changes},
Proc.\ R.\ Soc.\ Lond.\ A \textbf{392} (1984) 45.

\end{thebibliography}

\end{document}
